\documentclass[11pt,a4paper,oneside]{report}
\usepackage[T1]{fontenc}
\usepackage[utf8]{inputenc}
\usepackage[spanish,es-tabla]{babel}
\usepackage{lmodern}
\usepackage[a4paper,margin=2.5cm,headheight=15pt]{geometry}
\usepackage{microtype}
\usepackage{setspace}
\setstretch{1.05}
\usepackage{booktabs,tabularx,longtable,array}
\usepackage{xcolor}
\definecolor{academicblue}{HTML}{23445E}
\usepackage{fancyhdr}
\pagestyle{fancy}
\fancyhf{}
\fancyhead[L]{\small Actualidad en ciberseguridad}
\fancyfoot[C]{\thepage}
\renewcommand{\headrulewidth}{0.3pt}
\usepackage{csquotes}
\usepackage[backend=biber,style=numeric,sorting=none]{biblatex}
\addbibresource{referencias.bib}
\usepackage[unicode,hidelinks]{hyperref}
\hypersetup{pdftitle={Trabajo 1: Actualidad en ciberseguridad}}
\usepackage{bookmark}
\usepackage{titlesec}
\titleformat{\chapter}[hang]{\Large\bfseries\color{academicblue}}{\thechapter.}{0.6em}{}
\titlespacing*{\chapter}{0pt}{0pt}{14pt}
\titlespacing*{\section}{0pt}{9pt}{4pt}
\setcounter{tocdepth}{0}
% Activar solo si el profesor excluye el índice del límite de páginas.
\newif\ifmostrarindice
\mostrarindicefalse
% Anexos fuera del PDF de entrega por defecto.
\newif\ifmostraranexos
\mostraranexosfalse
\newcommand{\pendiente}[1]{\par\smallskip\noindent{\color{academicblue}\textbf{[PENDIENTE]} #1}\par\smallskip}
% Sustituir los campos antes de entregar.
\newcommand{\autor}{Enrique Sopeña}
\newcommand{\titulacion}{Máster Universitario en Ingeniería Informática}
\newcommand{\universidad}{Universitat Politècnica de València - ETSINF}
\newcommand{\curso}{20260 - 2027}
\newcommand{\fechaentrega}{25 de noviembre de 2026}
\newcommand{\asignatura}{Ciberseguridad}
\begin{document}
\begin{titlepage}
    \centering
    {\large\universidad\par}
    \vspace{0.4cm}
    {\large\titulacion\par}
    \vfill
    {\color{academicblue}\Huge\bfseries Actualidad en\\[0.2cm] ciberseguridad\par}
    \vspace{0.7cm}
    {\Large Trabajo 1\par}
    \vspace{0.5cm}
    {\large Análisis técnico y crítico de cinco incidentes\par}
    \vfill
    \begin{tabular}{rl}
        Autoría:    & \autor        \\
        Asignatura: & \asignatura   \\
        Curso:      & \curso        \\
        Entrega:    & \fechaentrega
    \end{tabular}
    \vspace{1.4cm}
\end{titlepage}
\setcounter{page}{2}
\ifmostrarindice
    \tableofcontents
    \clearpage
\fi
\chapter{Introducción y metodología}
\label{chap:metodologia}

\section{Objetivo y alcance}
La actualidad en ciberseguridad genera información
abundante, pero desigual en precisión y profundidad.
Comprender un incidente requiere distinguir el mecanismo
técnico, sus consecuencias y la evidencia que respalda
cada afirmación.

Este trabajo analiza cinco casos con actividad
documentada desde el 1 de septiembre de 2026:
PHP2Shell, Ivanti Connect Secure, los servicios
gubernamentales de Moldavia, la refrigeración comercial
en Francia y las entidades financieras surcoreanas.
El objetivo es responder a la siguiente pregunta:
¿qué debilidades, consecuencias y prioridades defensivas
se identifican al comparar estos incidentes?

Se examinan su desarrollo, impacto, mitigación y
relevancia, adaptando el análisis a la información
disponible. La investigación se cierra el
8 de noviembre de 2026. La fecha del incidente
se distingue de la publicación de noticias y avisos;
las vulnerabilidades anteriores se utilizan únicamente
como antecedentes de los ataques estudiados.
% Continúa el capítulo y la página de la introducción.
\section{Selección y verificación de fuentes}
La búsqueda combinó comunidades especializadas,
buscadores y publicaciones de organismos de seguridad.
Los debates se utilizaron para localizar candidatos;
las afirmaciones se contrastaron prioritariamente
con avisos de fabricantes, registros CVE, informes
de investigación y comunicados oficiales.

La selección atendió a cuatro criterios: actividad
documentada dentro del periodo, relevancia del impacto,
disponibilidad de evidencia y diversidad de mecanismos
y sectores. Se excluyeron noticias recientes que
describían incidentes anteriores al curso.
La muestra combina fallos de memoria, compromiso
de acceso remoto, extracción de datos y alteración
de procesos físicos; no pretende representar
estadísticamente todos los incidentes del periodo.

Para cada caso se reconstruyó la cronología y se
identificaron los sistemas afectados, el mecanismo
descrito y las medidas de respuesta.
Se contrastaron las publicaciones iniciales con
actualizaciones posteriores, diferenciando fuentes
complementarias de noticias que reproducían
un mismo comunicado.

\section{Criterios de análisis y limitaciones}
Se separaron hechos observados, valoraciones
de los investigadores e interpretaciones propias.
Las cifras se contextualizaron para no confundir
intentos, registros, personas y sistemas comprometidos.

Cuando existían CVE y CVSS publicados, se identificaron
su fuente y versión. Cuando faltaban, se valoró
la posibilidad de calcularlos: la estimación
de Moldavia se presenta como un escenario condicionado,
mientras que en los casos sin métricas suficientemente
documentadas se explicita la limitación.
Los valores de versiones distintas no se utilizan
para ordenar los incidentes.

El análisis del uso de IA distingue herramientas
observadas, funciones atribuidas y autonomía demostrada.
La comparación y las conclusiones se limitan
a la evidencia pública consultada y reconocen
la incertidumbre de investigaciones todavía abiertas.
\chapter{PHP2Shell: campaña de explotación de octubre}
\label{chap:incidente1}

\section{Contexto y justificación de selección}
Se analiza la campaña de octubre de 2026 que utilizó PHP2Shell,
dos vulnerabilidades del intérprete PHP, para comprometer
servidores y dispositivos de seguridad. Dream reconstruyó una
operación con 25 servidores comprometidos; Unit 42 confirmó las
herramientas y los resultados comunicados~\cite{dream2026php,unit422026php}.
Se selecciona por mostrar cómo un fallo en un componente compartido
puede convertir un dispositivo perimetral en una vía de acceso
a la infraestructura interna.

\section{Desarrollo del incidente y CVE/CVSS}
La operación combinó ejecución remota mediante PHP2Shell con
una escalada de privilegios del núcleo Linux,
CVE-2026-46300. Según Dream, los atacantes identificaron equipos
PAN-OS expuestos, seleccionaron versiones compatibles con su
explotación y comprometieron sus interfaces de gestión.
Los días 23 y 24 de octubre accedieron a dispositivos de un
entorno que los investigadores consideraron vinculado al
Gobierno estadounidense. Desde un cortafuegos establecieron
un túnel hacia la red interna y comprometieron un servidor
de gestión, obteniendo credenciales y enumerando
Active Directory~\cite{dream2026php}.

PHP2Shell comprende \textbf{CVE-2026-109618}, donde el agotamiento
de memoria deja datos de petición incorrectamente inicializados,
y \textbf{CVE-2026-109619}, donde otra petición atendida por el mismo
trabajador utiliza referencias a memoria liberada
(\textit{use-after-free}). La ejecución inicial tiene los permisos
del trabajador web; el acceso como administrador deriva del
encadenamiento con el fallo del núcleo~\cite{php2026oom,php2026stale,dream2026php}.

PHP Group publica para ambos \textbf{CVSS v4.0: 9,2, crítico},
con el vector base:\par
\noindent{\small\ttfamily
    CVSS:4.0/AV:N/AC:H/AT:P/PR:N/UI:N/\allowbreak
    VC:H/VI:H/VA:H/SC:N/SI:N/SA:N}\par
La valoración refleja acceso por red sin autenticación ni
interacción, complejidad alta y condiciones de ejecución
necesarias. No es un cálculo propio~\cite{php2026oom,php2026stale}.

\section{Impacto y mitigación}
Los compromisos permitieron acceso persistente, extracción de
credenciales y reconocimiento interno; no se acreditó el éxito
de todas las acciones posteriores~\cite{dream2026php}.
Para las organizaciones, esto compromete el control de sus
redes; para quienes utilizan sus servicios, supone riesgo de
exposición de información e interrupción del acceso.
No existe en las fuentes consultadas un balance global de
personas afectadas. Sí hubo consecuencias preventivas observables:
el NCSC neerlandés retiró temporalmente su web y comunicó
cierres similares en otras entidades~\cite{ncsc2026php}.

La corrección exige actualizar a PHP 8.2.35, 8.3.36, 8.4.27
o 8.5.12, según la rama, o aplicar los parches del proveedor.
Las modificaciones de configuración no ofrecen protección
completa~\cite{php2026aviso}. En PAN-OS deben instalarse sus
actualizaciones específicas y restringirse la gestión a redes
de confianza~\cite{unit422026php}.

\section{Análisis crítico}
Amazon observó intentos desde el 17 de octubre, sin confirmar
entonces su éxito~\cite{aws2026php}; CISA incorporó ambos CVE
al catálogo de vulnerabilidades explotadas el
día 26~\cite{cisa2026kev}.
Esta evolución obliga a distinguir exposición, intento y
compromiso. La principal lección es que parchear una
vulnerabilidad no elimina los accesos persistentes ni las
credenciales robadas: la recuperación debe investigar también
los sistemas alcanzados desde el dispositivo inicial.
\chapter{Ivanti: intrusión mediante una cadena de fallos}
\label{chap:incidente2}

\section{Contexto y justificación de selección}
El 5 de octubre de 2026, Volexity inició una investigación
que identificó actividad maliciosa desde el día 3 en una VPN
Ivanti Connect Secure 22.8R2.4, actualizada en aquel momento.
El atacante, denominado UTA0576, encadenó vulnerabilidades
para ejecutar comandos como administrador sin autenticarse.
Se selecciona porque muestra cómo comprometer una herramienta
de acceso remoto puede abrir el camino hacia los sistemas
que debería proteger~\cite{volexity2026ivanti}.

\section{Desarrollo del incidente y CVE/CVSS}
El atacante eludió los controles de autenticación y alcanzó
funciones internas de carga de archivos. Aprovechó el intervalo
entre la escritura y eliminación de un archivo rechazado para
ejecutar código. Después manipuló datos de configuración
utilizados por un servicio privilegiado y obtuvo permisos
\textit{root}~\cite{volexity2026ivanti}.

Los registros oficiales permiten distinguir tres componentes:
\textbf{CVE-2026-104935}, evasión de autenticación que permite
ejecutar comandos remotos, con \textbf{CVSS v3.1: 8,6};
\textbf{CVE-2026-118527}, nombres predecibles que permiten
depositar archivos en ubicaciones conocidas, con
\textbf{CVSS v3.1: 4,3}; y \textbf{CVE-2026-111357},
escalada de privilegios mediante confianza indebida en la
configuración DSCache, valorada con \textbf{CVSS v3.1: 7,0}
en el registro de MITRE. Son puntuaciones publicadas para
fallos individuales, no una valoración conjunta de la
cadena~\cite{cve2026ivanti104935,cve2026ivanti118527,
    cve2026ivanti111357}.

El primer fallo permite acceso por red, sin privilegios
previos ni interacción, según su vector oficial:\par
\noindent{\small\ttfamily
    CVSS:3.1/AV:N/AC:L/PR:N/UI:N/S:U/C:H/I:L/A:L}\par
\noindent La combinación elimina las barreras que limitarían
la explotación aislada de los otros
componentes~\cite{cve2026ivanti104935}.

\section{Impacto y mitigación}
Volexity documentó extracción de credenciales, exfiltración
de metadatos de Active Directory y búsquedas de contraseñas
administrativas y claves de recuperación BitLocker.
El atacante también borró registros y modificó componentes
en memoria; el comprobador de integridad de Ivanti no
detectaba esas modificaciones~\cite{volexity2026ivanti}.
Para las organizaciones afectadas, las credenciales robadas
pueden facilitar accesos posteriores y suplantaciones; para
sus usuarios, implican riesgos sobre información y servicios.
Las fuentes no ofrecen un recuento global de personas
afectadas.

La respuesta requiere instalar una versión que corrija
todos los componentes, según la rama: el aviso canadiense
actualizado recoge 25.1.1.3, 25.1.3.1 y 22.8R2.6.
CISA incorporó el primer CVE a su catálogo el 13 de octubre
y los otros dos el día 30~\cite{canada2026ivanti}.
Además de actualizar, deben investigarse los accesos
desde la VPN y revocarse las credenciales comprometidas;
una comprobación de integridad limpia no descarta la
intrusión~\cite{volexity2026ivanti}.

\section{Análisis crítico}
La principal lección es que estar actualizado no elimina
el riesgo de vulnerabilidades desconocidas. Tampoco basta
con valorar cada fallo por separado: una debilidad moderada
puede resultar decisiva dentro de una cadena.
La defensa debe combinar actualización, limitación de
privilegios y vigilancia de los accesos internos.
La ausencia de alteraciones detectables en disco debe
interpretarse junto con otras evidencias, no como garantía
de seguridad.
\chapter{Moldavia: extracción de datos gubernamentales}
\label{chap:incidente3}

\section{Contexto y justificación de selección}
Entre el 19 y el 23 de octubre de 2026, un ataque contra servicios públicos digitales de Moldavia permitió extraer datos personales. El incidente se comunicó públicamente el día 24~\cite{age2026inicial,age2026datos}. Se selecciona por sus consecuencias para la privacidad y la prestación de servicios públicos, y porque permite examinar los riesgos de conectar numerosos registros mediante aplicaciones con controles de acceso deficientes.

\section{Desarrollo del incidente y CVE/CVSS}
Según AGE, los atacantes explotaron servicios accesibles
sin autenticación y utilizaron la información obtenida
para realizar extracciones adicionales~\cite{age2026datos}.
La Presidencia describió fallos de autenticación y control
de acceso que permitieron consultar registros de al menos
diez instituciones. El ataque comenzó la tarde del 19;
el día siguiente se registró aproximadamente un millón
de solicitudes automatizadas, treinta veces el volumen
habitual~\cite{presidencia2026moldavia}.
Las fuentes no publican el código vulnerable ni permiten
identificar con certeza una técnica concreta, como
inyección SQL.

Los comunicados consultados no proporcionan CVE ni CVSS.
Como ejercicio propio, se estima un escenario de lectura
remota no autorizada siguiendo CVSS v3.1~\cite{firstcvss31}:\par
\noindent{\small\ttfamily
    CVSS:3.1/AV:N/AC:L/PR:N/UI:N/S:U/C:H/I:N/A:N}\par
\noindent El resultado es 7,5, alto, para este escenario hipotético de lectura no autorizada. Se asumen complejidad baja y alcance sin cambios; integridad y disponibilidad se fijan en cero para delimitar el cálculo, no porque se haya demostrado la ausencia de esos impactos. La documentación pública resulta insuficiente para asignar una puntuación definitiva a la vulnerabilidad real.

\section{Impacto y mitigación}
AGE estimó inicialmente extracciones asociadas a unos
750.000 identificadores personales. Entre las categorías
expuestas figuraban identificación, relaciones familiares,
prestaciones sociales e información financiera; no todas
afectaban a cada persona~\cite{age2026datos}.
Una evaluación posterior redujo el número de ciudadanos
afectados a aproximadamente 450.000, que debían ser
notificados~\cite{presidencia2026moldavia}.
La información facilita suplantaciones y fraudes dirigidos,
aunque conocer un IBAN no permite por sí solo disponer
del dinero. AGE no identificó modificaciones de los
registros ni exposición de datos de tarjetas
bancarias~\cite{age2026datos}.

La desconexión preventiva dejó servicios temporalmente
inaccesibles~\cite{age2026inicial}.
La respuesta posterior incluyó correcciones en MPass,
verificaciones adicionales de autorización y refuerzo
de las plataformas~\cite{presidencia2026moldavia}.
Como prevención se proponen controles de acceso por
operación y registro, límites de consultas y alertas
ante extracciones anómalas. Para los ciudadanos, AGE
recomendó verificar las solicitudes por canales oficiales
y no facilitar contraseñas ni códigos~\cite{age2026datos}.

\section{Análisis crítico}
El caso evidencia que contener el acceso no revierte
la pérdida de confidencialidad. La recuperación debe
combinar restauración segura, notificación y prevención
del fraude posterior.
También exige distinguir solicitudes, registros e individuos:
sus cifras no son equivalentes. La revisión del número
de afectados demuestra por qué deben fecharse las
estimaciones y actualizarse las conclusiones conforme
avanza la investigación.
\chapter{Francia: sabotaje de la refrigeración comercial}
\label{chap:incidente4}

\section{Contexto y justificación de selección}
Durante la noche del 11 al 12 de octubre de 2026,
varios cientos de supermercados franceses sufrieron
un ataque contra sus sistemas de refrigeración.
La activación de ciclos de descongelación alteró
la conservación de los productos y provocó pérdidas
de mercancía~\cite{onyphe2026francia}.
Se selecciona por mostrar cómo un incidente de
ciberseguridad puede afectar a procesos físicos,
a la distribución de alimentos y a la actividad
cotidiana de comercios y consumidores.

\section{Desarrollo del incidente y CVE/CVSS}
El CERT-FR confirmó compromisos de controladores
Danfoss AK-SM y AK-SC, además de sistemas Copeland
y Dixell XWEB~\cite{certfr2026frigorificos}.
Estos equipos permiten supervisar instalaciones y
modificar parámetros operativos desde una interfaz
web. En condiciones normales facilitan el mantenimiento
remoto; su compromiso permite que una acción digital
altere el funcionamiento de la
refrigeración~\cite{claroty2026danfoss}.

Existían vulnerabilidades documentadas en algunos
modelos. Claroty describió en el AK-SM 800A un mecanismo
alternativo de autenticación basado en un código
derivado de la fecha, así como fallos posteriores a
la autenticación que permitían ejecutar comandos
o modificar la configuración del servidor
web~\cite{claroty2026danfoss}.

Como antecedente técnico, \textbf{CVE-2025-41450}
afecta a la autenticación de la serie AK-SM 8xxA
anterior a la versión 4.2. Danfoss publica
\textbf{CVSS v3.1: 8,2, alto}, con el
vector~\cite{danfoss2025cve41450}:\par
\noindent{\small\ttfamily
    CVSS:3.1/AV:N/AC:H/PR:N/UI:R/S:C/C:H/I:L/A:H}\par
\noindent Esta puntuación corresponde al fallo conocido,
no al conjunto del incidente. Las fuentes consultadas
no permiten demostrar que fuera la vía utilizada
en todos los equipos atacados. Sin identificar
la vulnerabilidad concreta, no resulta defendible
asignar un CVSS propio al sabotaje.

\section{Impacto y mitigación}
La pérdida de refrigeración obligó a retirar alimentos,
con perjuicios económicos para los establecimientos
y desperdicio de productos~\cite{onyphe2026francia}.
Para los consumidores, la alteración de la cadena
de frío representa un riesgo sanitario si los productos
afectados llegan a consumirse; no se afirma que hubiera
intoxicaciones confirmadas.
El impacto combina manipulación de parámetros,
pérdida del servicio y consecuencias materiales.

El CERT-FR recomendó eliminar la exposición directa
a Internet, restablecer el acceso mediante una VPN
segura o una red protegida, actualizar el firmware
y cambiar las contraseñas predeterminadas.
Para equipos Danfoss compatibles, indicó utilizar
el nivel de seguridad \textit{Strict}.
Advirtió también que restaurar un equipo sin corregir
su exposición o vulnerabilidades no evita una nueva
intrusión, especialmente si ha terminado su
soporte~\cite{certfr2026frigorificos}.
Como prevención adicional, se propone supervisión
independiente de temperatura y alarmas ante cambios
anómalos de configuración.

\section{Análisis crítico}
La comodidad del mantenimiento remoto debe evaluarse
junto con las consecuencias de perder el control
del proceso físico. La protección requiere coordinar
al comercio, al mantenedor y al fabricante, definiendo
responsabilidades sobre accesos, actualizaciones
y respuesta.
Además, encontrar dispositivos expuestos o fallos
compatibles no demuestra el mecanismo exacto del
ataque: separar esas evidencias evita convertir
una hipótesis técnica en una causa confirmada.
\chapter{Corea del Sur: intrusiones financieras con IA}
\label{chap:incidente5}

\section{Contexto y justificación de selección}
CrowdStrike documentó una campaña contra entidades
financieras surcoreanas activa entre finales de septiembre
y principios de octubre de 2026. Su investigación,
publicada el 7 de octubre, relacionó la extracción de datos
con infraestructura que utilizaba ARTEX, una herramienta
de pruebas de penetración basada en agentes de
inteligencia artificial~\cite{crowdstrike2026artex}.
Se selecciona por aportar evidencia sobre el empleo
ofensivo de IA y permitir distinguir sus capacidades
observadas de las afirmaciones sobre ataques
completamente autónomos.

\section{Desarrollo del incidente y CVE/CVSS}
El análisis posterior de Zscaler ThreatLabz identifica
inyección SQL y evasión de autenticación en aplicaciones
web periféricas, incluidos portales de consulta de clientes
y servicios de aplicaciones móviles para empleados.
Según los investigadores, agentes basados en modelos
de lenguaje participaron en el reconocimiento y
la explotación~\cite{zscaler2026corea}.

La inyección SQL permite que entradas manipuladas
alteren las consultas a una base de datos; la evasión
de autenticación permite acceder sin superar correctamente
la comprobación de identidad. Son fallos de las aplicaciones
atacadas, no vulnerabilidades de ARTEX.
ThreatLabz describe una secuencia de reconocimiento,
explotación, extracción de información y búsqueda de
vías para comercializarla~\cite{zscaler2026corea}.

CrowdStrike aporta evidencias complementarias:
configuraciones de ARTEX e historiales de sesiones
de asistentes de programación encontrados en
infraestructura asociada al atacante.
Esto documenta el uso de IA, pero no demuestra que
todas las acciones fueran autónomas ni que cada
brecha del sector perteneciera a la misma
campaña~\cite{crowdstrike2026artex,fsc2026octubre}.

Los informes no identifican CVE ni publican CVSS
para los fallos concretos.
Se ha considerado una evaluación propia, pero conocer
la categoría del fallo no determina sus métricas:
faltan los privilegios necesarios, las condiciones
de explotación y los límites del sistema afectado.
No se asigna una cifra definitiva ni se utiliza
la cantidad de datos sustraídos como sustituto
de una evaluación técnica.

\section{Impacto y mitigación}
El 14 de octubre, la FSC contabilizaba brechas en
16 entidades y aproximadamente 100.000 registros
de clientes, con nombres, información de contacto
e identificadores personales. No se habían comunicado
filtraciones de contraseñas de cuentas ni números
de tarjetas~\cite{fsc2026octubre}.
El balance sectorial del 5 de noviembre ascendía
a 40 incidentes notificados en 37 entidades y unos
473.000 casos de exposición de información
personal~\cite{fsc2026noviembre}.
Estas cifras describen el conjunto sectorial, no
el alcance demostrado de la campaña de ARTEX.

Para los afectados, la exposición facilita
suplantaciones y fraudes dirigidos.
La FSC ordenó medidas de contención, parches,
suspensión de servicios vulnerables y protección
de clientes~\cite{fsc2026octubre}.
Posteriormente anunció inspecciones de API,
autenticación, controles de acceso y sistemas
periféricos, además de reforzar la vigilancia
de operaciones sospechosas~\cite{fsc2026noviembre}.
Como prevención adicional se propone limitar
privilegios y detectar consultas o extracciones
anómalas, independientemente de la herramienta
utilizada por el atacante.

\section{Análisis crítico}
La aportación del caso no es demostrar que la IA
pueda vulnerar cualquier banco, sino documentar
su integración en operaciones ofensivas reales.
La respuesta debe centrarse en los accesos y datos
que pueden comprometerse, sin depender de reconocer
una herramienta concreta.
Separar evidencia técnica, atribución y estadísticas
sectoriales evita exagerar la autonomía del atacante
y permite valorar las medidas defensivas con rigor.
\chapter{Comparación de los cinco casos}
\label{chap:comparativa}

\section{Síntesis de resultados}
La tabla~\ref{tab:comparativa} compara el mecanismo
documentado, las consecuencias y la prioridad defensiva.
Distingue los hechos observados de las limitaciones
de la información pública.

\begin{table}[htbp]
    \centering
    \small
    \renewcommand{\arraystretch}{1.15}
    \begin{tabularx}{\textwidth}{
        @{}p{2.1cm}XXX@{}
        }
        \toprule
        Caso                                  & Mecanismo & Impacto observado & Prioridad defensiva \\
        \midrule
        PHP2Shell                             &
        Fallos de memoria y encadenamiento con escalada
        de privilegios.                       &
        Compromiso de servidores, robo de credenciales
        y acceso interno.                     &
        Actualizar PHP y productos que lo integran;
        investigar accesos previos.                                                                 \\

        Ivanti                                &
        Evasión de autenticación, ejecución de código
        y escalada a \textit{root}.           &
        Extracción de credenciales y metadatos
        del directorio corporativo.           &
        Corregir toda la cadena, realizar análisis
        forense y revocar credenciales.                                                             \\

        Moldavia                              &
        Deficiencias de autenticación y autorización
        en servicios públicos.                &
        Extracción de datos personales e interrupciones
        preventivas.                          &
        Validar permisos por operación y registro;
        detectar consultas anómalas.                                                                \\

        Francia                               &
        Manipulación de controladores de refrigeración;
        vector concreto no acreditado.        &
        Alteración de la cadena de frío
        y pérdida de alimentos.               &
        Retirar la exposición directa y asegurar
        el mantenimiento remoto.                                                                    \\

        Corea del Sur                         &
        Inyección SQL y evasión de autenticación,
        según ThreatLabz, con apoyo de ARTEX. &
        Extracción de información personal
        y financiera.                         &
        Corregir aplicaciones periféricas y
        controlar accesos y extracciones.                                                           \\
        \bottomrule
    \end{tabularx}
    \caption{Comparación de los incidentes analizados.
        Síntesis basada en las fuentes de los capítulos anteriores.}
    \label{tab:comparativa}
\end{table}

\section{Patrones y diferencias}
El primer patrón es la importancia de las interfaces
accesibles desde el exterior. Las VPN, aplicaciones
periféricas y consolas de mantenimiento conectan
necesidades legítimas de acceso con activos sensibles.
Su exposición exige controles acordes con las
consecuencias de un compromiso, especialmente cuando
permiten alcanzar redes internas o modificar procesos
físicos~\cite{volexity2026ivanti,certfr2026frigorificos,
    zscaler2026corea}.

El segundo patrón es que contener el punto de entrada
no revierte todo el daño. En PHP2Shell e Ivanti,
la recuperación debe considerar credenciales y sistemas
alcanzados; en Moldavia y Corea del Sur, la información
extraída puede facilitar fraudes posteriores.
Francia añade una diferencia decisiva: restaurar
el controlador no recupera los alimentos
perdidos~\cite{dream2026php,volexity2026ivanti,
    age2026datos,zscaler2026corea,onyphe2026francia}.

Finalmente, los casos requieren respuestas distintas.
Los fallos conocidos demandan actualización y comprobación
de versiones; las deficiencias de autorización exigen
revisar qué puede consultar cada usuario.
El caso surcoreano incorpora IA a la ejecución ofensiva,
pero las debilidades descritas siguen requiriendo
controles de seguridad de aplicaciones
convencionales~\cite{zscaler2026corea}.

\section{Límites de la comparación}
No se establece una clasificación por CVSS:
PHP2Shell utiliza la versión 4.0 e Ivanti la 3.1,
mientras que otros casos carecen de puntuaciones
confirmadas para el fallo explotado.
Tampoco son equivalentes servidores comprometidos,
registros extraídos, personas afectadas y establecimientos.
La comparación es cualitativa y está limitada por
investigaciones abiertas y niveles desiguales
de detalle técnico.
\chapter{Conclusiones}
\label{chap:conclusiones}

\section{Respuesta al objetivo}
El análisis de cinco incidentes ocurridos desde septiembre
de 2026 muestra que la relevancia de un suceso de
ciberseguridad depende tanto del fallo técnico como
del contexto en el que se explota. Sus consecuencias
abarcan el acceso a redes internas, la exposición
de información personal, la interrupción de servicios
y la pérdida de bienes físicos.

PHP2Shell e Ivanti ilustran cómo el encadenamiento
de vulnerabilidades puede ampliar los permisos
y el alcance de una intrusión. Moldavia y las entidades
financieras surcoreanas evidencian la importancia
de controlar el acceso a datos desde aplicaciones
y servicios periféricos. Los supermercados franceses
muestran que la seguridad de los sistemas conectados
también protege procesos materiales y la conservación
de alimentos.

Los casos permiten concluir que actualizar es necesario,
pero no suficiente para gestionar un incidente.
La corrección del punto de entrada debe acompañarse
de una investigación del alcance, la recuperación
de los sistemas y medidas frente a las consecuencias
que persisten, como credenciales sustraídas
o datos personales ya extraídos.

\section{Recomendaciones y reflexión final}
A partir de la comparación se proponen tres prioridades.
En primer lugar, las organizaciones deben conocer
sus activos y reducir la exposición de interfaces
administrativas, aplicaciones auxiliares y equipos
sin soporte. En segundo lugar, deben aplicar
actualizaciones y verificar los permisos efectivos
sobre funciones y registros, evitando que un acceso
limitado permita consultar información ajena
o alcanzar sistemas adicionales. En tercer lugar,
la respuesta debe integrar evidencias técnicas,
protección de los afectados y comunicación clara
sobre lo confirmado y lo todavía incierto.

El empleo de IA observado en la campaña surcoreana
añade capacidad de automatización, pero no sustituye
la necesidad de corregir las debilidades explotadas.
Su presencia debe analizarse mediante evidencias,
sin equiparar automáticamente asistencia,
automatización y autonomía completa.

Finalmente, la investigación ha mostrado que
el rigor también consiste en reconocer los límites
de las fuentes. Un CVSS no mide el daño social,
un sistema expuesto no equivale a una víctima
y una hipótesis compatible no demuestra la causa
de un ataque. Las conclusiones se limitan a los
cinco casos estudiados y a la información pública
consultada hasta el 8 de noviembre de 2026;
no constituyen una descripción exhaustiva
del panorama de amenazas.
\clearpage
% Mostrar un marcador si aún no existen referencias citadas.
\defbibheading{bibliography}[Bibliografía]{\chapter*{#1}\addcontentsline{toc}{chapter}{#1}}
\chapter*{Bibliografía}\addcontentsline{toc}{chapter}{Bibliografía}
\printbibliography[heading=none]
\pendiente{Incorporar referencias verificadas a referencias.bib y citarlas en el texto.}
\ifmostraranexos
    \appendix
    \chapter{Registro de búsqueda y selección}
\pendiente{Registrar por candidato: título, fecha del hecho, fecha de publicación, fuente, criterio de inclusión o descarte y justificación.}
\chapter{Trazabilidad de evidencias}
\pendiente{Registrar por afirmación: caso, afirmación, clave bibliográfica, localización exacta, fecha de consulta, corroboración y grado de certeza. Incluir solo material necesario y respetar la atribución.}
\chapter{Justificación de las evaluaciones CVSS}
\pendiente{Para cada evaluación propia: versión, herramienta utilizada, vector, métricas, evidencia, supuestos, cálculo reproducible y sensibilidad ante datos inciertos.}
\chapter{Material complementario}
\pendiente{Añadir cronologías, diagramas propios, glosario y evidencias adicionales pertinentes, siempre referenciados desde el cuerpo principal.}

\fi
\end{document}
